\chapter{Oncologist}
\index{Oncologist}

\medskip
\noindent\textit{\textbf{Under review.} This chapter is an AI-assisted educational draft; it carries no source citations and is not a substitute for professional medical advice.}

\section*{Definition and Overview}
An oncologist is a medical specialist who diagnoses, treats, and manages patients with cancer. In many countries, oncology is a specialist medical area recognised by the Royal Australasian College of Physicians (RACP), with three main subspecialty streams: medical oncology, radiation oncology, and surgical oncology. Medical oncologists deliver systemic therapies (chemotherapy, targeted therapy, immunotherapy, hormonal therapy); radiation oncologists prescribe and oversee radiotherapy; surgical oncologists perform cancer surgery. The oncologist coordinates the multidisciplinary team (MDT) that plans each patient's care using tumour board discussion, pathology, staging, and evidence-based guidelines. Beyond active treatment, oncologists provide survivorship care, symptom and palliative management, and participation in clinical trials to advance treatment.

\section*{Epidemiology}
Cancer is a leading cause of death in many countries, with more than 160,000 new cases and over 50,000 cancer deaths annually. Approximately one in two people is expected to be diagnosed with cancer by age 85, and cancer care is delivered by several thousand practising oncologists across some countries. Medical oncology and radiation oncology have grown substantially with the ageing population and advances in systemic therapy. Access to oncology care varies geographically, with rural and remote people facing reduced access to specialist centres and rural outreach/heli-medical services. The number of practising oncologists continues to rise, and the field increasingly integrates genomics, precision medicine, and immunotherapy into routine care. Multidisciplinary tumour boards are standard across local cancer services.

\section*{Aetiology and Pathophysiology}
Oncology is the management of malignant disease rather than a single pathological process, but its subspecialty structure reflects mechanism of treatment:
\begin{itemize}
    \item \textbf{Medical oncology:} systemic drug therapy -- cytotoxic chemotherapy, targeted agents (e.g., tyrosine kinase inhibitors such as imatinib, osimertinib), immunotherapy (checkpoint inhibitors such as pembrolizumab, nivolumab), hormonal agents (e.g., tamoxifen, anastrozole), and antibody--drug conjugates
    \item \textbf{Radiation oncology:} uses ionising radiation to kill cancer cells via DNA damage; includes external-beam radiotherapy, brachytherapy, stereotactic radiosurgery; delivered with planning (CT-based dosimetry)
    \item \textbf{Surgical oncology:} resection of primary tumours and metastases with curative intent (e.g., lumpectomy/mastectomy, colectomy, Whipple procedure), and cytoreduction
    \item \textbf{Supportive and palliative care integration:} symptom control, pain management, and end-of-life care alongside active treatment
    \item \textbf{Precision oncology:} tumour genomic profiling and biomarker testing (e.g., HER2, EGFR, BRAF, PD-L1, MSI) to guide targeted and immunotherapies
    \item \textbf{Clinical trials and research:} all local oncology centres participate in trials to develop new treatments
\end{itemize}

\section*{Clinical Features}
The ``clinical picture'' of an oncologist is the scope of professional care:
\begin{itemize}
    \item \textbf{Diagnosis and staging:} reviewing histopathology, imaging, and biomarker results; assigning stage to guide treatment
    \item \textbf{Treatment planning:} recommending surgery, radiotherapy, systemic therapy, or combinations; discussing options, benefits, and risks with the patient
    \item \textbf{Systemic therapy administration:} chemotherapy, immunotherapy, targeted and hormonal therapy; managing infusions and side-effects
    \item \textbf{Toxicity management:} neutropenia and infection, nausea/vomiting, neuropathy, rash, immune-related adverse events, cytopenias
    \item \textbf{Surveillance and follow-up:} scheduled imaging and blood tests to detect recurrence after treatment
    \item \textbf{Survivorship:} managing long-term effects, cancer-related fatigue, and surveillance for second cancers
    \item \textbf{MDT coordination:} leading or participating in multidisciplinary tumour board meetings with surgeons, pathologists, radiologists, and nurses
\end{itemize}

\section*{Differential Diagnosis}
Considerations in differential diagnosis relevant to oncology practice:
\begin{itemize}
    \item \textbf{Benign vs malignant:} distinguishing benign tumours and non-malignant mimics (e.g., sarcoidosis, tuberculosis, amyloid) from cancer
    \item \textbf{Subspecialty pathways:} medical vs radiation vs surgical oncology vs haemato-oncology (haematologist managing leukaemia/lymphoma)
    \item \textbf{Palliative care physician:} focuses on symptom management and quality of life rather than anti-cancer therapy
    \item \textbf{General physician/hospitalist:} provides general medical care, not cancer-specific therapy
    \item \textbf{Oncologist vs oncologist-in-training:} fellows working toward RACP specialist recognition
\end{itemize}

\section*{When to Seek Medical Care}
Patients undergoing oncology treatment should seek care promptly for new or worsening symptoms, fevers, bleeding, or treatment side-effects. Those with a new cancer diagnosis should be referred to an oncologist through their GP or surgeon for multidisciplinary assessment.

\textbf{Contact your local emergency services for:}
\begin{itemize}
    \item Fever or rigors during chemotherapy (possible neutropenic sepsis) -- immediate
    \item Severe bleeding, uncontrolled vomiting, or dehydration
    \item Chest pain, breathlessness, or neurological symptoms (possible PE, stroke, or cord compression)
    \item Anaphylaxis or severe infusion reaction
    \item Signs of spinal cord compression (new back pain with leg weakness or urinary retention)
\end{itemize}

\section*{Diagnosis and Investigations}
\begin{itemize}
    \item \textbf{Pathology:} histology and molecular/genomic profiling of the tumour
    \item \textbf{Imaging:} CT, PET-CT, MRI, and bone scans for staging and response assessment
    \item \textbf{Tumour markers:} disease- and tumour-specific markers (e.g., PSA, CA125, CEA) used in monitoring
    \item \textbf{Staging systems:} TNM staging and clinical guidelines (e.g., eviQ) to plan treatment
    \item \textbf{Baseline organ function:} FBC, renal and liver function, cardiac assessment (echocardiogram) before certain cardiotoxic therapies
    \item \textbf{Biomarker testing:} PD-L1, HER2, EGFR, BRAF, BRCA, MSI/MMR status to select targeted and immunotherapy
    \item \textbf{Clinical trial screening:} assessing eligibility for investigational protocols
\end{itemize}

\section*{Management}
\begin{itemize}
    \item \textbf{Multidisciplinary care:} tumour board planning integrating surgery, radiotherapy, and systemic therapy
    \item \textbf{Medical oncology:} neoadjuvant, adjuvant, and palliative chemotherapy; targeted agents; immunotherapy; hormonal therapy -- regimens per eviQ protocols
    \item \textbf{Radiation oncology:} curative or palliative radiotherapy; e.g., whole-breast radiotherapy after lumpectomy, stereotactic radiotherapy for brain metastases
    \item \textbf{Surgical oncology:} resection with clear margins; sentinel node biopsy; cytoreductive surgery
    \item \textbf{Toxicity prophylaxis:} growth factors (G-CSF), antiemetics (ondansetron, aprepitant), antimicrobial prophylaxis in high-risk regimens
    \item \textbf{Symptom management:} analgesia, antiemetics, laxatives, management of neuropathy and fatigue
    \item \textbf{Survivorship and surveillance:} scheduled imaging, bloods, and long-term follow-up; screening for treatment-related second cancers
    \item \textbf{Palliative care referral:} for advanced disease to optimise comfort and quality of life alongside or instead of active therapy
    \item \textbf{Clinical trials:} participation strongly supported within local research centres
\end{itemize}

\section*{Complications}
\begin{itemize}
    \item Neutropenic sepsis and opportunistic infection during chemotherapy
    \item Cytopenias (anaemia, thrombocytopenia) requiring transfusion or growth factors
    \item Organ toxicity: cardiotoxicity (anthracyclines, trastuzumab), nephrotoxicity (cisplatin), hepatotoxicity, pneumonitis
    \item Immune-related adverse events from checkpoint inhibitors (colitis, hepatitis, endocrinopathies, pneumonitis)
    \item Peripheral neuropathy (e.g., platinum agents, taxanes), mucositis, nausea, and severe fatigue
    \item Treatment-related second primary cancers and chemotherapy-induced infertility
    \item Radiation-related late effects (fibrosis, second malignancy, organ injury)
\end{itemize}

\section*{Prognosis and Outcomes}
Prognosis in oncology is tumour-specific and stage-dependent. Overall five-year relative survival for all cancers in many countries now exceeds 70\%, up from around 50\% in the mid-1980s, reflecting earlier detection and improved treatment. Survival varies widely by cancer type and stage: prostate and melanoma survival is very high, while pancreatic and lung cancer survival lags. Immune checkpoint inhibitors and targeted therapies have markedly improved outcomes in melanoma, lung, kidney, and other cancers. The oncologist's role in coordinating evidence-based, multidisciplinary care and clinical-trial access is central to these gains. Ongoing research continues to improve cure rates and quality of life for people living with cancer.

\section*{Prevention and Self-Care}
\begin{itemize}
    \item Maintain healthy lifestyle: non-smoking, limiting alcohol, sun protection, physical activity, healthy body weight
    \item Participate in population screening (bowel, breast, cervical) and discuss individualised risk with a GP
    \item Attend scheduled oncology follow-up and surveillance appointments
    \item Report new symptoms or treatment side-effects early to the treating team
    \item Maintain vaccination status (influenza, pneumococcal, COVID-19) as advised during treatment
    \item Engage supportive services: psychology, dietetics, physiotherapy, cancer community organisations (Cancer Council the affected region)
    \item Consider advance care planning and discuss preferences with family and the oncology team
\end{itemize}

\section*{Sources and Further Reading}
The following reputable sources provide further information for healthcare students and patients seeking reliable, current advice about oncology in Australia.

\begin{itemize}
    \item Cancer Council Australia. \textit{Cancer information, support, and research.} https://www.cancer.org.au/
    \item eviQ Cancer Treatments Online. \textit{Evidence-based cancer treatment protocols.} https://www.eviq.org.au/
    \item Royal Australasian College of Physicians (RACP). \textit{Medical oncology training.} https://www.racp.edu.au/
    \item Royal Australian and New Zealand College of Radiologists (RANZCR). \textit{Radiation oncology.} https://www.ranzcr.com/
    \item Clinical Oncology Society of Australia (COSA). \textit{Professional oncology resources.} https://www.cosa.org.au/
    \item Australian Institute of Health and Welfare (AIHW). \textit{Cancer data and statistics.} https://www.aihw.gov.au/
    \item Healthdirect Australia. \textit{Cancer and oncology services.} https://www.healthdirect.gov.au/
\end{itemize}

\clearpage
